%%%%%%%%%%%%%%%%%%%%%%%%%%%%%%%%%%%%%%%%%%%%%%%%%%%%%%%%%%%%%%%%%%%%%%
% Lớp 12 | Phương trình mặt phẳng, đường thẳng, mặt cầu | Phương trình mặt cầu trong không gian
% Dạng 2H5?3-1: Câu hỏi lý thuyết
% Số câu: 12  (sinh tự động bởi tach_cau.py; từ nay thầy sửa trực tiếp file này)
%%%%%%%%%%%%%%%%%%%%%%%%%%%%%%%%%%%%%%%%%%%%%%%%%%%%%%%%%%%%%%%%%%%%%%

%%%%%%%%%%%%%%%% TRẮC NGHIỆM 4 PHƯƠNG ÁN %%%%%%%%%%%%%%%%

% nguồn: Phan1_Lop_12_Hinh_hoc_Chuong_5
\begin{ex}%[2H5N3-1]%[ma:a909b431]
	Trong không gian $Oxyz$, cho mặt cầu $(S)$ có tâm~là điểm $I(a; b; c)$ và bán kính  $R$. Phương trình của $(S)$~là
	 \choice
		{\True $(x-a)^2 + (y-b)^2 +  (z-c)^2 =R^2$}
		{$(x+a)^2 + (y+b)^2 + (z+c)^2  =R^2$}
		{$(x-a)^2 + (y-b)^2  + (z-c)^2 =R$}
		{$(x+a)^2 + (y+b)^2  + (z+c)^2 =R$}
	 \loigiai{
		Mặt cầu $(S)$ có tâm $I(a; b;  c)$, bán kính $R$ có phương trình~là $$(x-a)^2 + (y-b)^2 + (z-c)^2 =R^2.$$
	 }
\end{ex}

% nguồn: Phan1_Lop_12_Hinh_hoc_Chuong_5
\begin{ex}%[2H5N3-1]%[ma:e3e94c5a]
	Trong không gian $Oxyz$, phương trình nào trong các phương trình sau  là phương trình của một mặt cầu?
	 \choice
		{$(S)\colon x^2+y^2-z^2-4x+2y+2z-3=0$} 
		{$(S)\colon  x^2+2y^2+z^2-4x+2y+2z-10=0$}
		{\True $(S)\colon  x^2+y^2+z^2-4x+2y+2z+2=0$}
		 {$(S)\colon x^2+y^2+z^2-4x+2y+2z+8=0$}
	\loigiai{ 
		Phương trình $x^2+y^2+z^2-2ax-2by-2cz+d=0$ $(a^2+b^2+c^2-d>0)$ là phương trình mặt  cầu có tâm $I(a;b;c)$, bán kính $R=\sqrt{a^2+b^2+c^2-d}$.\\
		Trong phương trình $(S)\colon x^2+y^2+z^2-4x+2y+2z+2=0$, ta có $\heva{&a=2\\&b=-1\\&c=-1\\&d=2} \Rightarrow a^2+b^2+c^2-\mathrm{d}=4>0$ nên  phương trình này là phương trình mặt cầu.
	} 
\end{ex}

% nguồn: Phan1_Lop_12_Hinh_hoc_Chuong_5
\begin{ex}%[2H5N3-1]%[ma:cfdb2939]
	Trong  không gian với hệ toạ độ $Oxyz$, phương trình nào sau đây là phương trình của một mặt cầu?
	 \choice
		{$-(x-1)^2+(y+1)^2+(z-10)^2=225$} 
		{$(x-1)^2-(y+1)^2+(z-10)^2=225$} 
		 {$(x-1)^2+(y+1)^2-(z-10)^2=225$}
		 {\True $(x-1)^2+(y+1)^2+(z-10)^2=225$}
	\loigiai{ 
		Phương  trình mặt cầu là $(x-1)^2+(y+1)^2+(z-10)^2=225$.
	 }
\end{ex}

% nguồn: Phan1_Lop_12_Hinh_hoc_Chuong_5
\begin{ex}%[2H5N3-1]%[ma:db3003fb]
	Trong không  gian $Oxyz$, phương trình nào trong các phương trình sau là phương trình tổng quát của một mặt phẳng?
	\choice 
		 {$\dfrac{1}{x}+\dfrac{2}{y}+\dfrac{3}{z}=1$}
		 {$\dfrac{x}{2}=\dfrac{y-1}{4}=z$}
		 {$x^2+y^2+z^2-1=0$}
		 {\True $2x-3z+1=0$}
	\loigiai{Phương trình  $2x-3z+1=0$ là phương trình tổng quát của một mặt phẳng}
\end{ex}

% nguồn: Phan1_Lop_12_Hinh_hoc_Chuong_5
\begin{ex}%[2H5N3-1]%[ma:50e75ed5]
	Trong không gian $Oxyz$, phương trình $x^2 +y^2 +z^2-2ax-2by-2cz+d=0$~là phương trình của  một mặt cầu khi và chỉ khi
	 \choice
		{\True $a^2  +b^2 +c^2 -d>0$}
		{$a^2 +b^2 +c^2 +d  \ge 0$}
		 {$a^2 +b^2 +c^2 +d>0$}
		{$a^2 +b^2 +c^2 -d  \ge 0$}
	 \loigiai{
		Phương trình $x^2 +y^2 +z^2-2ax-2by-2cz+d=0$~là phương trình của một mặt cầu khi và chỉ khi $a^2 +b^2  +c^2 -d>0$.
	 }
\end{ex}

% nguồn: Phan1_Lop_12_Hinh_hoc_Chuong_5
\begin{ex}%[2H5H3-1]%[ma:498d60d7]
	Trong không gian $Oxyz$, có tất  cả bao nhiêu giá trị nguyên của tham số $m$ để \break $x^2+y^2+z^2+2(m+2)x-2(m-1)z+3m^2-5=0$ là phương trình của một mặt cầu?
	 \choice
		{$5$} 
		{\True $7$} 
		 {$6$}
		{$4$} 
	 \loigiai{
		Ta có $a=-m-2$, $b=0$, $c=m-1$,  $d=3m^2-5$.\\
		$a^2+b^2+c^2-d=(-m-2)^2+(m-1)^2-3m^2+5=-m^2+2m+10$.\\ 
		Để phương trình $x^2+y^2+z^2+2(m+2)x-2(m-1)z+3m^2-5=0$ là phương trình mặt  cầu thì $$-m^2+2m+10>0\Leftrightarrow 1-\sqrt{11}<m<\sqrt{11}+1.$$
		 Vì $m\in\mathbb{Z}$ nên $ m\in\{-2;-1;0;1;2;3;4\}$.
	} 
\end{ex}

% nguồn: Phan1_Lop_12_Hinh_hoc_Chuong_5
\begin{ex}%[2H5H3-1]%[ma:02df30d6]
	Trong không gian $Oxyz$, cho mặt cầu $(S)\colon\left(x-2\right)^2+\left(y+3\right)^2+\left(z-1\right)^2=9$. Điểm nào trong các điểm  bên dưới thuộc mặt cầu $(S)$?
	\choice 
		 {\True $K\left(5;-3;1\right)$}
		 {$J\left(-2;3;-1\right)$}
		{$H\left(-7;-3;1\right)$} 
		 {$I\left(2;-3;1\right)$}
	\loigiai{ 
		Thay tọa độ các điểm vào phương trình mặt cầu $(S)$  ta có\\
		Điểm $K\left(5;-3;1\right)\in(S)$ vì $\left(5-2\right)^2+\left(-3+3\right)^2+\left(1-1\right)^2=9$  $\Leftrightarrow $ $9=9$ (luôn đúng).\\
		Điểm  $J\left(-2;3;-1\right)\notin(S)$ vì $\left(-2-2\right)^2+\left(3+3\right)^2+\left(-1-1\right)^2=9$ $\Leftrightarrow $ $56=9$ (vô lý).\\
		Điểm $H\left(-7;-3;1\right)\notin(S)$ vì $\left(-7-1\right)^2+\left(-3-2\right)^2+\left(1-3\right)^2=9$ $\Leftrightarrow $ $93=9$ (vô lý).\\ 
	Điểm $I\left(2;-3;1\right)\notin(S)$ vì  $\left(2-1\right)^2+\left(-3-2\right)^2+\left(1-3\right)^2=9$ $\Leftrightarrow $ $30=9$ (vô lý).}
\end{ex}

% nguồn: Phan1_Lop_12_Hinh_hoc_Chuong_5
\begin{ex}%[2H5H3-1]%[ma:ca04dd7d]
	Trong không gian $Oxyz$, mặt cầu  $(S)\colon\left(x-2\right)^2+\left(y-1\right)^2+\left(z+3\right)^2=16$ đi qua điểm nào dưới đây?
	 \choice
		{$Q\left(-2;-1;-1\right)$} 
		{$N\left(-2;-1;3\right)$} 
		 {$M\left(2;1;-3\right)$}
		{\True $P\left(2;1;1\right)$} 
	\loigiai{ 
	Thay tọa độ điểm $P\left(2;1;1\right)$ vào phương trình mặt cầu $(S)\colon\left(x-2\right)^2+\left(y-1\right)^2+\left(z+3\right)^2=16$ (thỏa mãn). Ta có mặt cầu $(S)$  đi qua điểm $P$.}
\end{ex}

% nguồn: Phan1_Lop_12_Hinh_hoc_Chuong_5
\begin{ex}%[2H5H3-1]%[ma:42d78a2a]
	Trong không gian $Oxyz$, cho mặt cầu $(S)\colon\left(x-2\right)^2+\left(y+1\right)^2+\left(z-1\right)^2=6$. Điểm nào dưới đây thuộc  mặt cầu $(S)$?
	 \choice
		{$B\left(3;1;1\right)$} 
		 {$A\left(3;-2;2\right)$}
		 {\True $C\left(3;-2;3\right)$}
		 {$D\left(1;0;4\right)$}
	 \loigiai{
	 Ta thấy $\left(3-2\right)^2+\left(-2+1\right)^2+\left(3-1\right)^2=6\Rightarrow C\left(3;-2;3\right)\in(S)$.}
\end{ex}

% nguồn: Phan1_Lop_12_Hinh_hoc_Chuong_5
\begin{ex}%[2H5H3-1]%[ma:4c55cbc0]
	Trong không gian  $Oxyz$, phương trình nào sau đây là phương trình mặt cầu?
	\choice 
		{$x^2+y^2+z^2-4y-2z+5=0$} 
		{\True  $x^2+y^2+z^2-4x+1=0$}
		 {$x^2+y^2+z^2-6z+20=0$}
		{$x^2+y^2+z^2+2x+6y-2z+15=0$} 
	\loigiai{ 
		Xét phương trình $x^2+y^2+z^2-4x+1=0$, ta  có $a=2$, $b=c=0$ và $d=1$.\\
		Khi đó,  ta có
		\[a^2+b^2+c^2-d=2^2+0^2+0^2-1=3  >0.\]
	Vậy phương trình $x^2+y^2+z^2-4x+1=0$ là  phương trình mặt cầu tâm $I(2;0;0)$ và $R=\sqrt{3}$.}
\end{ex}

% nguồn: Phan1_Lop_12_Hinh_hoc_Chuong_5
\begin{ex}%[2H5H3-1]%[ma:b60ed42e]
	Trong không gian với hệ trục tọa độ $Oxyz$, cho hai  mặt cầu $(S)\colon\left(x-3\right)^2+y^2+z^2=9$ và $\left(S'\right)\colon\left(x+2\right)^2+y^2+z^2=4$. Khẳng định nào sau đây là đúng?
	\choice 
		{\True Hai mặt cầu  tiếp xúc ngoài}
		 {Hai mặt cầu tiếp xúc trong}
		{Hai  mặt cầu không có điểm chung}
		{Hai mặt cầu có nhiều hơn  một điểm chung}
	 \loigiai{
		Mặt cầu $(S)$ có tâm $I\left(3;0;0\right)$, $R=3$ còn mặt cầu $\left(S'\right)$ có  tâm $I'\left(-2;0;0\right)$, $R'=2$.\\
	Do $II'=5=R+R'$ nên hai  mặt cầu tiếp xúc ngoài.}
\end{ex}

% nguồn: Phan1_Lop_12_Hinh_hoc_Chuong_5
\begin{ex}%[2H5H3-1]%[ma:19d2d38a]
	Trong không  gian $Oxyz$, phương trình nào dưới đây là phương trình của một mặt cầu?
	 \choice
		{\True $x^2+y^2+z^2-2 x-2 y-2  z-1=0$}
		{$x^2+y^2+z^2-2 x-2  y-2 z+3=0$}
		{$x^2+y^2-z^2-2 x-2 y-2 z-1=0$} 
		{$x^2+y^2+2 z^2-2 x-2 y-2  z-1=0$}
	 \loigiai{
		Phương  trình $x^2+y^2+z^2-2 x-2 y-2 z-1=0$ thỏa mãn $a^2+b^2+c^2-d=1^2+1^2+1^2+1=4>0$ nên đây là phương trình mặt cầu.
	 }
\end{ex}
